%!TEX TS-program = xelatex
%!TEX encoding = UTF-8 Unicode
% Awesome CV LaTeX Template for CV/Resume
%
% Derived from resume.tex for a PhD application to the Medical Data Science
% Group at ETH Zurich (Prof. Julia Vogt). Do not edit resume.tex from here.
%


%-------------------------------------------------------------------------------
% CONFIGURATIONS
%-------------------------------------------------------------------------------
\documentclass[11pt, a4paper]{../awesome-cv}
\usepackage{../cv-common}

% Body copy is Roboto Thin (\bodyfontlight), which reads as pale gray.
% Same sizes and layout. Regular weight in near-black for body text only.
\makeatletter
\renewcommand*{\paragraphstyle}{\fontsize{8pt}{1em}\bodyfont\upshape\color{black}}
\renewcommand*{\descriptionstyle}[1]{{\fontsize{8pt}{1em}\bodyfont\upshape\color{black} #1}}
\renewcommand*{\skillsetstyle}[1]{{\fontsize{8pt}{1em}\bodyfont\color{black} #1}}
\makeatother

%-------------------------------------------------------------------------------
%	PERSONAL INFORMATION
%-------------------------------------------------------------------------------
\name{Jaume}{Ivars Grimalt}
\position{Machine Learning Research | Computer Vision | Medical Imaging}
% Location + relocation intent. No freelance / part-time availability.
\address{Currently in Valencia, Spain | EU citizen | Open to relocation to Zurich, Munich, London}

\mobile{+34 690 183 644}
\email{jaumeig95@gmail.com}
\homepage{jaumeivars.com}
\github{jaume2000}
\linkedin{jaume-ivars-grimalt}
\googlescholar{b5gMLtMAAAAJ}{Google Scholar}

\photo[circle,edge,left]{../PFP.jpeg}


%-------------------------------------------------------------------------------
\begin{document}

% Print the header with above personal information
\makecvheader[C]

% Print the footer with 3 arguments(<left>, <center>, <right>)
\makecvfooter
  {\today}
  {Jaume Ivars Grimalt~~~·~~~CV}
  {\thepage}


%-------------------------------------------------------------------------------
%	CONTENT
%-------------------------------------------------------------------------------
\cvsection{Profile}

\begin{cvparagraph}
Computer vision engineer moving into research. My current work studies representational dynamics in residual vision networks: how features evolve with depth and what that reveals about weight sharing and model internals. I have applied deep learning to medical and biological data, from brain MRI segmentation at MIALAB (UPV) to a microbiology detection system I built from scratch and tested in five laboratories.
\end{cvparagraph}

\cvsection{Education}
\nopagebreak
\begin{cventries}
\cvkeep
\cvdegree
    {UPV · GPA 8.7/10 · Thesis Honours (became Mycrospace)}
    {M.Sc. Artificial Intelligence, Pattern Recognition \& Digital Imaging}
    {Valencia, Spain}
    {2023 – 2024}
% TODO: thesis title, supervisor
\cvkeep
\cvdegree
    {CDTM (TUM/LMU) · 16 selected}
    {Technology Management}
    {Valencia / Munich}
    {2024}
% TODO: thesis title, supervisor
\cvkeep
\cvdegree
    {UPV · GPA 8.3/10 · Honours in ML and Statistics}
    {B.Sc. Computer Engineering}
    {Valencia, Spain}
    {2019 – 2023}
% TODO: thesis title, supervisor
\end{cventries}

\cvsection{Research}
\par\nopagebreak
\vspace{\acvSectionContentTopSkip}
\vspace{-2.5mm}
{\bodyfont\fontsize{7.6pt}{0.95em}\selectfont\color{black}\raggedright
\noindent\textbullet~~\href{https://jaumeivars.com/research/representation-dynamics}{\textit{What Extra Euler Steps Reveal About Residual Networks: Feature Dynamics in Weight-Tied and Untied Stages.}}
Manuscript in preparation. \href{https://jaumeivars.com/research/representation-dynamics}{jaumeivars.com/research/representation-dynamics}.
Studies continuous-depth dynamics in ConvNeXt, Swin and ResNet: finer numerical integration without retraining, trajectory geometry and per-block weight corrections.
With \textasciitilde\textbf{9.6M} fewer parameters, a weight-shared ConvNeXt-T stage trained from scratch on ImageNet-1k reaches \textbf{80.34\%} top-1 (baseline: \textbf{80.96\%}).\par
}
% TODO: arXiv link

\cvsection{Publications}
\par\nopagebreak
\vspace{\acvSectionContentTopSkip}
\vspace{-2.5mm}
{\bodyfont\fontsize{7.6pt}{0.95em}\selectfont\color{black}\raggedright
\noindent\textbullet~~\textit{Auto-N2N: Self-supervised MRI Denoising using Synthetic Pair Generation and Mixture of Experts.}
S.\ Morell-Ortega et al., incl.\ \textbf{J.\ Ivars Grimalt}. Preprint, 2026. \href{https://openreview.net/forum?noteId=1f93N1n8RP}{OpenReview}\par
\vspace{0.15mm}
\noindent\textbullet~~\textit{Detecting Cytopathic Effect in Cell Culture with Computer Vision.}
W.\ Ragazzo, \textbf{J.\ Ivars Grimalt}. Under review.\par
}

\cvsection{Experience}
\begin{cventries}
\cvkeep
\cventry
    {Co-Founder \& CTO}
    {Mycrospace}
    {Valencia, Spain}
    {2024 – Oct 2026}
    {
\begin{cvitems}
\item {Started as a Master's thesis in 2024. The company was founded in January 2026. Computer vision for microbiology, tested in \textbf{5 labs}. 1st Prize IdeasUPV (2024).}
\item {\textbf{Model.} Own ConvNeXt detector (YOLO/DQ-DETR-inspired), from scratch; ~10–2000 colonies/image. \textbf{67 mAP50-95}; 92\% recall / 93\% precision.}
\item {\textbf{Training.} \textbf{12,500+} GPU-hours on Leonardo (two EuroHPC AI Factory allocations). \textbf{4x A100}, DDP/AMP/torch.compile/SLURM, checkpoint-resumed jobs. Own framework: k-fold CV, WBF ensembles.}
\item {\textbf{Data.} \textbf{1,500}/\textbf{5,000} images annotated with microbiologists. Weakly-supervised labelling (click $\rightarrow$ crop $\rightarrow$ 4-ch CNN heatmap $\rightarrow$ bbox); versioning + k-fold label scoring.}
\item {\textbf{Product/infra.} BullMQ/Redis inference on GCP, plus NestJS/PostgreSQL, Next.js and React Native.}
\end{cvitems}
    }
\cvkeep
\cventry
    {Research Engineer (full-time)}
    {MIALAB: Medical Imaging Analysis Laboratory, UPV}
    {Valencia, Spain}
    {May 2025 – Jan 2026}
    {
\begin{cvitems}
\item {Volumetric brain MRI segmentation with nnU-Net; multi-scale feature strategies. Profiled volBrain (\textbf{700k+} volumes) from NIfTI ingest to voxel prediction for inference cost.}
\end{cvitems}
    }
\end{cventries}
\vspace{0.4mm}
\noindent{\fontsize{8pt}{1em}\bodyfont\bfseries\color{darktext} Earlier experience:}~~{\fontsize{8pt}{1em}\bodyfont\color{black} Tapstar, Full-Stack Engineer (freelance), 2025; Vesta-Z, Data Engineer, 2023–2025; Neurocatching, Data Scientist Intern, 2023; Sciling, Backend Developer Intern, 2022.}\par

\cvsection{Technical Skills}
\begin{cvskills}
  \cvskill
    {AI \& ML}
    {Object detection, instance segmentation, density estimation, classification, CNNs, ViTs, foundation models, Neural ODEs, PINNs}
  \cvskill
    {Medical imaging}
    {nnU-Net, volumetric MRI, NIfTI}
  \cvskill
    {Training}
    {Python, PyTorch (DDP, AMP, torch.compile), Hugging Face, timm, multi-GPU HPC (SLURM)}
  \cvskill
    {Production}
    {ONNX, Docker, CI/CD}
\end{cvskills}

\cvsection{Languages}
\begin{cvparagraph}
Spanish (native), Catalan (native), English (C1)
\end{cvparagraph}

%-------------------------------------------------------------------------------
\end{document}
